Synthetic data generation for healthcare is motivated by strict privacy regulations, data scarcity for rare conditions, and the need for high-fidelity statistical structure in complex clinical datasets.
Existing generative approaches---including GANs, VAEs, and neural diffusion models---often struggle with the heterogeneous, mixed-type, and constraint-rich nature of Electronic Health Records (EHRs), while requiring substantial GPU resources and architectural tuning.

This project develops a two-stage, tree-based generative framework for synthesising realistic Intensive Care Unit (ICU) patient trajectories from the MIMIC-III clinical database. The system combines flow matching with gradient-boosted trees (XGBoost) to learn a time-indexed vector field that transforms noise into tabular data, avoiding the smoothness assumptions that hinder neural generators on mixed-type features. Two model backbones are implemented: ForestFlow, which treats all target dimensions as continuous flows, and Heterogeneous Sequential Feature Forest Flow (HS3F), which generates features sequentially using flow-matching regressors for continuous variables and XGBoost classifiers for categorical variables.

Generation follows an autoregressive factorisation: an IID hour-0 model synthesises initial patient states, and a conditional model generates subsequent hourly transitions given lagged dynamics and static covariates. The framework is evaluated on three axes: downstream utility via multi-task Train-on-Synthetic, Test-on-Real (TSTR) evaluation covering patient-level mortality and ICU length-of-stay prediction; distributional fidelity via Kolmogorov--Smirnov statistics, correlation preservation, and clinical range violation rates; and privacy risk via Distance-to-Closest-Record diagnostics and DOMIAS-style membership inference attacks.

Empirical evidence comes from two protocols on MIMIC-III: a single-seed HS3F sweep over a $5 \times 8$ grid ($n_t \in \{1,3,5,7,9\}$, $n_{\text{noise}} \in \{1,3,5,7,9,11,13,15\}$) with $35$ cells currently complete, and a three-seed matched-cell comparison of HS3F, ForestFlow, and a CTGAN/CRGAN lower-bound baseline at $(n_t{=}5, n_{\text{noise}}{=}5)$. Time-level discretisation $n_t$ is the dominant hyperparameter: the best single-seed configuration reaches mortality ROC-AUC $\approx 0.74$ (versus $\approx 0.83$ real-model baseline), a utility gap of approximately $0.09$, while clinical range violations drop from $\approx 30\%$ at low $n_t$ to $\approx 14\%$ at $n_t \geq 7$. Macro-F1 remains low ($\leq 0.30$) across all cells, indicating minority-class collapse under ICU class imbalance. The matched cell favours HS3F on ranking utility (ROC-AUC $0.67$) and on trajectory-level temporal coherence (autocorrelation distance $\approx 2\times$ lower than ForestFlow), while CTGAN collapses to chance-level ranking (ROC-AUC $\approx 0.50$). Membership-inference ROC-AUC remains near chance throughout ($0.48$--$0.50$) with zero exact-match synthetic records, and CPU-only training completes in minutes to a few hours per configuration. Extrapolation analysis suggests the remaining utility gap is structural rather than hyperparameter-limited. These findings highlight both the promise and current limitations of tree-based flow matching for sequential clinical data synthesis.
